% !TeX root = ../../main.tex
\subsection{Guide dependence limits gene-level calibration}
\label{sec:calibration}

The preceding benchmarks use control diagnostics at different reporting
levels.  To examine the effect of aggregation directly, we simulated all-null
beta-binomial screens across sample sizes, mean guide counts, and dispersions,
with three or five guides per gene.  A focused arm introduced within-gene
guide correlation of 0.4.  Odd-numbered null genes estimated a control scale
after Stouffer aggregation, and even-numbered null genes evaluated it.  The
full grid and a separate continuous-dose diagnostic are described in
Supplementary Section~\ref{sec:supp-calibration}.

Under independent guides, the largest five-seed mean gene-level $p<0.05$ rate
was 0.074.  In the correlated-guide arm, guide-level rejection rates remained
near or below 0.05, while gene-level rates reached 0.108--0.262
(Table~\ref{tab:calibration-grid}).  Aggregation-matched held-out scaling
reduced these rates to 0.056--0.112 but did not remove the inflation,
particularly with five guides per gene.  These results locate a source of
miscalibration in the independence assumption of the gene combiner.  A
near-nominal guide-level diagnostic alone is insufficient to establish
calibration of the reported gene probabilities.

\begin{table}[htbp]
\centering
\caption{Null calibration before and after gene aggregation.  Entries are
ranges of five-seed mean rejection rates across the applicable grid cells.
The held-out column applies a scale learned from a disjoint half of the null
genes after the same aggregation.  ``Mod. denotes guide-variance
moderation.  Rejection uses $p<0.05$, with a target rate of 0.05.}
\label{tab:calibration-grid}
\small
\setlength{\tabcolsep}{3.5pt}
\begin{tabular}{llrccc}
\toprule
Guide dependence & Guides & Mod. & Guide & Gene & Held-out gene\\
\midrule
Independent & 3 & No  & 0.000--0.058 & 0.002--0.064 & 0.000--0.076\\
Independent & 3 & Yes & 0.000--0.060 & 0.002--0.074 & 0.004--0.088\\
Independent & 5 & No  & 0.000--0.056 & 0.004--0.064 & 0.004--0.064\\
Independent & 5 & Yes & 0.000--0.060 & 0.026--0.074 & 0.012--0.060\\
Correlated ($r=0.4$) & 3 & No  & 0.027--0.043 & 0.124--0.144 & 0.060\\
Correlated ($r=0.4$) & 3 & Yes & 0.000--0.030 & 0.108--0.140 & 0.056--0.068\\
Correlated ($r=0.4$) & 5 & No  & 0.029--0.048 & 0.232--0.254 & 0.104--0.112\\
Correlated ($r=0.4$) & 5 & Yes & 0.000--0.041 & 0.224--0.262 & 0.068--0.076\\
\bottomrule
\end{tabular}
\end{table}

\paragraph{Denominator sensitivity in an external audit.}
An audit of the benchmark deposited by \citet{dempster2025fdr,dempster2025fdrdata}
examined a separate source of error: changing library totals by restricting
the input to null genes before fitting.  Because H.-H.J.\ also developed
CB\textsuperscript{2}, this is an audit of an author's prior method.
For the same Avana null genes, restoring full-library CB\textsuperscript{2}
probabilities reduced discoveries at FDR 0.10 from 152 to zero.  BARCS was not
fitted in that comparison.  The largest cell-line-specific null $p<0.05$ rate
remained 0.145, so the corrected discovery count does not establish nominal
calibration.

BARCS analyses of the deposited normalized PSN1 and DeWeirdt tables provided
mixed diagnostics.  PSN1 produced no gene calls and guide-level $p<0.05$ rates
of 0.0481 before scaling and 0.0465 after scaling.  In DeWeirdt, BARCS recovered
one of 28 curated interactions, and 37 of its 38 calls arose in one comparison.
The complete audit is reported in Supplementary
Section~\ref{sec:supp-fdr-audit}.  Together with the null grid, these results
show why denominator preservation, guide dependence, and independent control
evaluation must each be addressed when interpreting a discovery threshold.
